% Источники: exam/materials/моделирование.pdf, разделы 14--17.
\section{Классификация математических моделей.}

В материалах курса явно выделены следующие группы математических моделей:
статические и динамические, вероятностные и статистические, аналитические и
имитационные, стохастические и марковские модели.

\textbf{Аналитическая модель} -- математическая модель объекта или процесса
реального мира, записанная с помощью аналитических выражений с учетом
физико-химических взаимодействий в исследуемой системе.

\textbf{Динамическая модель} описывает, как протекает явление или изменяется
объект во времени, то есть в динамике. При изучении динамических моделей
задается начальное состояние объекта или процесса, затем исследуются
изменения.

Динамические модели обладают следующими свойствами:
\begin{enumerate}
  \item временная зависимость: модели учитывают изменения параметров
  системы во времени;
  \item начальные условия: начальные состояния системы определяют
  дальнейшее ее поведение;
  \item устойчивость: определение, сохраняется ли поведение системы при малых
  возмущениях.
\end{enumerate}

\textbf{Статическая модель} отражает состояние исследуемой системы или объекта
в заданный момент времени.

\textbf{Статистические модели} предполагают решение основных задач
математической статистики:
\begin{itemize}
  \item получение точечных оценок некоторой случайно переменной;
  \item построение доверительных интервалов для случайных параметров;
  \item проверка статистических гипотез.
\end{itemize}

\textbf{Вероятностные модели} предполагают решение задачи теории вероятности.
Задачи математической статистики являются обратными по отношению к задачам
теории вероятности: в теории вероятности после задания некоторого случайного
явления требуется рассчитать вероятностные характеристики этого явления или
процесса.

Если система многокомпонентна, то есть описывается взаимодействием ее
составных частей, аналитические модели оказываются слишком трудоемкими и
ресурсозатратными при их реализации на ЦВМ. Альтернативой является
\textbf{имитационное моделирование}, не учитывающее конкретных
физико-химических процессов, протекающих в системе, а лишь приближающее
закон изменения входных данных при получении определенного набора выходных
данных.

\textbf{Марковская модель} -- это таблица или граф переходов марковской цепи с
описанием каждого из состояний системы и указанием вероятностей перехода
между состояниями.

\clearpage
